%%%%%%%%%%%%%%%%%%%%%%%%%%%%%%%%%%%%%%%%%%%%%%%%%%%%
\section{Related Work}

\subsection{Frequent Itemset Mining}

Apriori~\cite{agrawal1994fast} and FP-Growth~\cite{han2000mining}
are foundational methods for mining frequent itemsets.
ECLAT~\cite{zaki1997new} offers a vertical data format alternative.
A comprehensive survey of FIM is given in~\cite{fournier2017survey}.
All of these methods evaluate support over the entire dataset and
cannot detect itemsets that occur at high frequency only within
specific time intervals.

\subsection{Dense Interval Mining}

Methods that output patterns together with their dense intervals
fall into two categories.
\textit{Occurrence-gap-based} methods include
LPFIM~\cite{yin2014lpfim}, LPPM~\cite{saleh2011lppm}, and
RPM~\cite{kiran2017rpm}.
LPFIM partitions intervals when occurrence gaps reach a threshold;
LPPM introduces cumulative spillover for noise-robust detection;
RPM groups consecutive occurrences with gaps within a maximum period.
These methods indirectly estimate density through gap constraints
or periodicity thresholds, making parameter tuning difficult and
results inexact with respect to the direct sliding-window frequency.
\textit{Direct} approaches include GLFMiner~\cite{fournierviger2021glf},
which mines patterns at fixed temporal granularities and thus cannot
identify intervals of arbitrary width, and SIM~\cite{kiran2015sim},
which returns only the single optimal interval per itemset.
None of these methods support streaming or provide a mechanism
to control scan frequency.

\subsection{Stream Mining}

Sliding-window-based stream mining has been studied in
incremental FIM~\cite{lee2001sliding,leung2006dstree},
where windows limit the temporal scope of support computation.
However, these methods still report global support within the window
and do not enumerate all dense sub-intervals.

